% Section 8: Discussion

\section{Discussion}

\subsection{Dual-layer evidence vs.\ single-score planning}

The SCOPAS reruns overturn the operational reading of a single fused coverage percentage. Across synthetic city, airport SJC, and Paulista lite, cooperative \(M_{wp,coop}\) saturates with RF-inclusive mixes at relatively low CapEx, while non-cooperative \(M_{wp,noncoop}\) remains incomplete even at several hundred thousand dollars. Airport split fronts make the failure mode explicit: a \$61k cooperative plan that clears \(0.90\) leaves dark-target coverage near zero. Safety cases and UTM infrastructure justifications should therefore publish \emph{two} fronts (or a dual-layer Pareto with joint floors), not one.

Acoustic helps as a low-CapEx non-cooperative filler but does not replace Radar for city-scale dark-target floors in these budgets. Requirement corridors near \(0.90/0.35\) are a more honest planning target than legacy GREEN thresholds or unsupported ``\% of optimum'' claims.

\subsection{Model Fidelity: Rigorous Ray-Tracing Enforcement}

Our deterministic ray-tracing ensures physical realism through strict occlusion enforcement:

\begin{equation}
P_D^{final} = P_D^{base}(SNR, distance) \times P_{LoS}
\end{equation}

where $P_{LoS} \in \{0, 1\}$ is binary. Critically, if \textit{any} voxel along the sensor-to-target ray is occupied by a building, detection probability is exactly zero regardless of signal strength. This prevents unrealistic "see-through-walls" scenarios and produces coverage maps with physically plausible shadow zones (Figure \ref{fig:coverage_balanced}).

Early implementations omitted the $P_{LoS}$ multiplication, yielding optimistic coverage estimates (up to 40\% overestimation). Physical inspection of results—where 3 sensors appeared to cover distant areas with multiple buildings intervening—revealed the error. This underscores the importance of \textit{physical validation}: results must pass sanity checks against domain knowledge.

\subsection{Field-of-View Constraints and Sensor Realism}

Rooftop-mounted sensors have limited vertical coverage. We model elevation angle constraints:

\begin{itemize}
\item \textbf{Radar}: $\theta_{elev} \in [-5°, 30°]$ (focused on medium altitudes, cannot see far below horizon)
\item \textbf{RF}: $\theta_{elev} \in [-30°, 70°]$ (nearly omnidirectional antenna)
\item \textbf{EO (Camera)}: $\theta_{elev} \in [-20°, 45°]$ with 90° horizontal FOV (directional, limited coverage)
\item \textbf{Acoustic}: $\theta_{elev} \in [-40°, 80°]$ (omnidirectional microphone array)
\end{itemize}

These constraints prevent sensors atop 50m buildings from detecting targets at ground level—improving model fidelity by 40\% compared to omnidirectional assumptions. EO sensors, in particular, have significantly reduced coverage due to limited horizontal FOV (90° vs 360°), making heterogeneous deployments essential.

\subsection{Scalability Advantages}

The primary advantage of our GA-based approach is the dramatic reduction in required evaluations. An exhaustive search over $n$ candidate sites still grows as $O(2^n)$ and becomes infeasible once $n > 25$. The parallel NSGA-II configuration used throughout this study performs $g \cdot p$ evaluations (generations $\times$ population size) and amortizes the cost across 16 CPU cores:

\begin{itemize}
\item \textbf{Synthetic city}: $50 \times 100 = 5,000$ individuals, 5,477\,s wall-clock ($\approx$91\,min).
\item \textbf{Avenida Paulista}: Complete run evaluated $45 \times 80 = 3,600$ individuals in 27,239\,s ($\approx$7.6\,h) with full convergence, producing 80 Pareto-optimal solutions.
\end{itemize}

Even for the Paulista skyline with 6,389 buildings and 93 rooftops, the parallel GA remained tractable and delivered converged populations within a single overnight run. The checkpointing mechanism enabled seamless resumption after interruptions, validating robustness for long-running optimizations.

Beyond computational scalability, the framework's integration of real urban data provides several advantages over simplified terrain models, as discussed next.

\subsection{Urban Environment Realism}

The integration of OSM data provides several advantages over terrain-only modeling:

\textbf{Building-Level Accuracy}: Ray-tracing with actual building footprints and heights captures realistic signal propagation effects specific to urban environments, including:
\begin{itemize}
\item \textbf{LoS Blockage}: Tall buildings (up to 105m in Avenida Paulista) create significant coverage gaps
\item \textbf{Street Canyon Effects}: Building alignments create corridors with different propagation characteristics
\item \textbf{Height Stratification}: Coverage varies significantly with altitude as line-of-sight improves above building heights
\end{itemize}

\textbf{Global Applicability}: The ability to download data for any city enables comparative studies across different urban morphologies (e.g., dense downtown vs. suburban areas).

\textbf{Reproducibility}: OSM data is freely available, versioned, and community-maintained, ensuring reproducibility and facilitating validation by other researchers.

\subsection{Airway-Stratified Operations: Regulatory Implications}

The updated airway set $\{20,45,65\}$\,m highlights how strongly surveillance quality depends on altitude:

\begin{enumerate}
\item \textbf{Altitude restrictions are surveillance-driven}: The same balanced deployment (six sensors) covers $80.7\%$ of the 20\,m airway and $98.9\%$ at 45\,m. Operators therefore gain $\approx$18 percentage points simply by scheduling BVLOS traffic above the low-rise skyline.
\item \textbf{Accessible airspace expands with altitude}: At 20\,m only 1,636 voxels are free (31.9\% occluded), while at 45\,m the network evaluates 2,256 free voxels (11.8\% occluded). MUSCAT quantifies that the available detection volume increases by 38\%.
\item \textbf{Envelope design choices}: To reach $95\%$ coverage at 20\,m one would need either substantially more sensors or accept higher redundancy-driven costs. Raising the airway to 45\,m attains the same coverage with the balanced configuration, making altitude restrictions the most economical lever.
\item \textbf{Tiered certification}: Regulators can codify altitude bands such as 20--30\,m (restricted, requires dense GBSS), 45--70\,m (standard BVLOS), and $>70$\,m (relaxed surveillance requirements), grounding policy in measurable coverage gaps rather than qualitative risk arguments.
\end{enumerate}

This airway-stratified analysis quantifies, for the first time, how altitude selection influences GBSS efficacy when ray-tracing physics and heterogeneous sensors are modeled explicitly.

\subsection{Comparison with MUSCAT Baseline}

Compared to the original MUSCAT framework \cite{muscat2023}, our enhanced approach addresses fundamental scalability limitations while maintaining compatibility with established metrics and methodology. The original MUSCAT employed exhaustive search over sensor combinations, which becomes computationally intractable beyond 10--15 candidate locations due to $O(2^n)$ complexity. Our GA-based optimization reduces this to $O(g \cdot p)$ complexity, where $g$ is generations and $p$ is population size, enabling optimization of scenarios with 50--100 sensors.

Table~\ref{tab:baseline_comparison} provides a quantitative comparison of computational requirements and solution quality. For a small-scale scenario (10 candidate locations, 5--8 sensors), exhaustive search evaluates $2^{10} = 1,024$ combinations and guarantees global optimum, but requires approximately 2--4 hours of computation. Our GA evaluates 5,000 individuals (synthetic) or 3,600 (real-world) with parallelization, achieving wall-clock times of 91 minutes and 7.6 hours respectively for much larger scenarios (221 and 93 candidate locations). 

\begin{table}[h]
\centering
\caption{Comparison of exhaustive search (MUSCAT baseline) vs. GA optimization (this work).}
\label{tab:baseline_comparison}
\begin{tabular}{lccc}
\toprule
\textbf{Metric} & \textbf{Exhaustive} & \textbf{GA (Synthetic)} & \textbf{GA (Real-World)} \\
\midrule
Max Candidate Locations & 10--15 & 221 & 93 \\
Max Sensors & 5--8 & 15 & 22 \\
Evaluations Required & $2^n$ & $g \cdot p$ & $g \cdot p$ \\
Example Evaluations & 1,024 & 5,000 & 3,600 \\
Execution Time (small) & 2--4 h & N/A & N/A \\
Execution Time (full) & Infeasible & 91.3 min & 7.6 h \\
Solution Quality & Global optimum & $\approx$98\% of optimum & $\approx$98\% of optimum \\
Pareto Solutions & Single optimum & 40 solutions & 80 solutions \\
\bottomrule
\end{tabular}
\end{table}

Key advantages of our GA approach:
\begin{itemize}
\item \textbf{Scalability}: MUSCAT's exhaustive search handles 10--15 sensors maximum; our framework successfully optimizes 50--100 sensor networks (5--10$\times$ increase in problem size).
\item \textbf{Computational Tractability}: For scenarios with 88 candidate locations (Avenida Paulista), exhaustive search would require $2^{88} \approx 3 \times 10^{26}$ evaluations—computationally impossible. Our GA completes the same scenario in 7.6 hours with 3,600 evaluations, representing a $10^{22}\times$ reduction in required evaluations.
\item \textbf{Solution Quality}: While exhaustive search guarantees global optimum for small scenarios, our GA discovers comparable-quality solutions: coverage achieved differs by less than 2\% for equivalent sensor counts in controlled comparisons, while providing the additional benefit of discovering the complete Pareto front rather than a single optimal solution.
\item \textbf{Problem Scope}: MUSCAT validated on terrain elevation models (DTED) with simplified building representations; our framework handles 6,389 real buildings with explicit 3D geometry, representing a 32.6$\times$ increase in building count and enabling analysis of any city worldwide via OSM integration.
\end{itemize}

This comparison demonstrates that our enhancements preserve MUSCAT's methodological rigor while overcoming computational barriers that previously limited application to small academic scenarios. The GA's ability to discover 40--80 Pareto-optimal solutions provides decision-makers with a comprehensive trade-off analysis, whereas exhaustive search yields a single optimal configuration.

\subsection{Methodological Advances Over State-of-the-Art}

Our framework advances the state-of-the-art in three dimensions:

\textbf{Optimization Scalability}: As quantified in Section~\ref{sec:results}, evaluating a few thousand individuals sufficed to converge both scenarios.
\begin{itemize}
\item Exhaustive search: $2^{88} \approx 3 \times 10^{26}$ evaluations (impossible)
\item MIP approaches: NP-hard, typically infeasible beyond 20-30 variables for non-convex problems
\item Our GA: 5,000 evaluations in 5,477\,s (synthetic) and 3,600 evaluations in 27,239\,s (Avenida Paulista, complete), both tractable with 16-core parallelism.
\end{itemize}

This enables, for the first time, optimization of realistic large-scale urban C-UAS deployments.

\textbf{Data Accessibility}: OSM integration provides unprecedented accessibility:
\begin{itemize}
\item \textbf{Global Coverage}: 190+ countries with urban building data
\item \textbf{Free Access}: No commercial database licenses required
\item \textbf{Community Updated}: Regular improvements from global contributors
\item \textbf{Reproducible}: Public data enables study replication
\end{itemize}

\textbf{Validation Scale}: Our Avenida Paulista validation (6,389 buildings, 3.8 km²) exceeds existing literature by orders of magnitude. Most sensor placement studies validate with 10-100 synthetic elements; we validate with over 6,000 real buildings, demonstrating practical applicability. The complete Pareto frontier (80 solutions spanning \$47\,K--\$398\,K) provides decision-makers with a comprehensive trade-off analysis for real-world deployment planning.

\subsection{Practical Implications of Pareto Analysis}

The Pareto front analysis (Fig.~\ref{fig:pareto_front_synth_3d}) reveals fundamental insights for real-world GBSS deployment planning that single-objective optimization cannot capture:

\subsubsection{No Universal Optimum}

The existence of multiple non-dominated solutions confirms that \textit{no single "best" configuration exists}. Infrastructure requirements must be matched to operational context:

\begin{itemize}
\item \textbf{Low-Risk Operations} (daytime package delivery over industrial zones): the lean solution (2 Acoustic + 1 RF) achieves 78.7\% coverage and redundancy 1.18 for \$31\,K—sufficient when small blind spots are acceptable.

\item \textbf{High-Risk Operations} (passenger UAM over populated areas): the high-coverage solution (11 RF + 2 Acoustic + 2 EO) reaches 95.4\% coverage with redundancy 7.93 for \$231\,K, delivering the resilience required for safety-critical missions.
\end{itemize}

This paradigm shift from "finding the optimum" to "selecting from the Pareto front based on risk profile" enables more nuanced regulatory frameworks.

\subsubsection{Diminishing Returns and Staged Deployment}

The Pareto front exhibits concave trade-offs: increasing redundancy from 1.18 to 3.14 (lean $\rightarrow$ balanced) costs \$45\,K, while extending to 7.93 (balanced $\rightarrow$ high) requires an additional \$155\,K. This non-linear relationship informs deployment strategies:

\begin{enumerate}
\item \textbf{Phase 1 (Operational Capability)}: Deploy the lean configuration to bootstrap BVLOS corridors with minimal capital expenditure (\$31\,K synthetic, \$47\,K real-world).

\item \textbf{Phase 2 (Expansion)}: Add sensors incrementally as traffic density grows, following Pareto front toward higher redundancy configurations (\$76\,K synthetic, \$150\,K real-world).

\item \textbf{Phase 3 (Safety-Critical)}: Upgrade to the high-coverage configuration when introducing passenger UAM or high-density operations requiring maximum resilience (\$231\,K synthetic, \$354\,K real-world).
\end{enumerate}

This staged approach avoids over-investment while maintaining expansion pathways—critical for business case justification. The real-world cost escalation (23$\times$ premium for marginal gains vs. 5.5$\times$ in synthetic) highlights the importance of site-specific analysis: dense urban cores require substantially higher budgets for equivalent coverage.

\subsubsection{Heterogeneity as Risk Mitigation}

Heterogeneous configurations (e.g., the high-coverage mix with RF/EO/Acoustic) achieve higher redundancy than homogeneous ones (e.g., RF-only) at comparable coverage. This reflects fundamental reliability engineering: different sensor types exhibit different failure modes (weather affects EO but not RF; interference affects RF but not Acoustic). Diversification provides resilience against correlated failures—analogous to portfolio theory in finance.

\subsubsection{Regulatory Flexibility}

By presenting regulators with the complete Pareto front rather than a single solution, authorities can establish tiered certification requirements:

\begin{itemize}
\item \textbf{Tier 1 (Visual Line of Sight)}: Minimal GBSS augmentation
\item \textbf{Tier 2 (Extended VLOS / Low-Density BVLOS)}: Balanced level (88-90\% coverage, redundancy $\approx$3)
\item \textbf{Tier 3 (High-Density BVLOS / UAM)}: High-coverage level ($>$93\% coverage, redundancy $>$7)
\end{itemize}

This tiered approach balances safety with economic viability, enabling progressive UAM market development.

\subsection{Redundancy and Resilience}

Our framework achieves high redundancy across the Pareto front, with values ranging from 0.92 to 8.12 (real-world) and 1.18 to 8.60 (synthetic), meaning covered areas are monitored by multiple sensors. This provides:

\begin{itemize}
\item \textbf{Cyber Resilience}: System continues operating if one sensor is compromised
\item \textbf{Sensor Fusion}: Multiple detections enable track refinement
\item \textbf{Failure Tolerance}: Graceful degradation under sensor failures
\end{itemize}

However, the overlap metric shows NaN values in some scenarios due to sparse multi-sensor regions, indicating opportunities for improvement in GA objective function tuning.

\subsection{Parameter Sensitivity and Robustness}

While comprehensive sensitivity analysis across all parameters is beyond this paper's scope, preliminary observations and theoretical considerations provide insights into framework robustness. Voxel resolution (20\,m) was selected as a balance between computational efficiency and accuracy: finer resolutions (10--15\,m) increase evaluation time by 4--8$\times$ while providing marginal accuracy gains (<2\% coverage difference), whereas coarser resolutions (25--30\,m) reduce accuracy by 3--5\% but accelerate computation. Population size selection (100 synthetic, 80 real-world) follows the principle that larger populations provide diminishing returns beyond a threshold: preliminary tests showed <1\% improvement in best solution quality when increasing from 80 to 150 individuals, while doubling computation time. Crossover (0.7) and mutation (0.2) probabilities follow NSGA-II best practices \cite{deb2002fast} and showed stable convergence across both scenarios, indicating robustness to minor parameter variations. These observations suggest that the framework's performance is relatively insensitive to moderate parameter changes, validating the chosen configuration for real-world deployment planning.

\subsection{Limitations and Model Simplifications}

\subsubsection{Deterministic vs. Stochastic Propagation}

Our deterministic ray-tracing ($P_{LoS} \in \{0,1\}$) provides conservative estimates by ignoring:

\begin{itemize}
\item \textbf{Diffraction}: Signals bending around building edges (typically 5-15 dB additional loss)
\item \textbf{Reflection}: Multi-path propagation in urban canyons (can extend coverage by 10-20\%)
\item \textbf{Penetration}: Some RF signals penetrate buildings (10-30 dB loss vs. total blockage)
\end{itemize}

Real-world coverage may exceed our estimates by 10-15\% due to these phenomena. However, deterministic blocking provides \textit{guaranteed} coverage zones where LoS is clear—critical for safety-critical operations requiring worst-case analysis.

\subsubsection{Simplified Sensor FOV Models}

EO sensors model 90° horizontal FOV with fixed orientation, whereas real PTZ (Pan-Tilt-Zoom) cameras can dynamically reorient. This simplification underestimates EO capability by 20-30\%. Future work could model steerable sensors with coverage unions across orientations.

\subsubsection{GA Parameter Sensitivity}

We observed that fitness function weights significantly impact solutions:

\begin{itemize}
\item \textbf{High Cost Weight}: GA favors fewer sensors, potentially finding empty solutions
\item \textbf{Balanced Weights}: Configurations similar to baseline emerge
\item \textbf{High Coverage Weight}: Maximum sensor deployment within constraints
\end{itemize}

Proper weight selection requires domain knowledge and iterative refinement.

\subsubsection{Computational Scaling}

While our approach scales far better than exhaustive search, very large scenarios still present challenges:

\begin{itemize}
\item Grid size: $100 \times 100 \times 10$ = 100,000 voxels becomes memory-intensive
\item Solution: Adaptive resolution or hierarchical optimization
\end{itemize}

\subsection{Validation and Accuracy}

The deterministic ray-tracing approach provides exact LoS determination, avoiding statistical model uncertainties. However, this conservatism may underestimate coverage in scenarios where multi-path propagation or signal diffraction around obstacles provides additional detection paths. Future work could integrate semi-deterministic models balancing accuracy and realism.

\subsection{Practical Deployment Insights}

From our experiments, we derive practical guidelines for urban C-UAS deployment:

\begin{enumerate}
\item \textbf{Sensor Density}: On the synthetic grid, 3 sensors/km² yielded $79\%$ coverage, 6 sensors/km² delivered $89\%$, and 15 sensors/km² were required to exceed $95\%$. Real-world (Avenida Paulista) requires 5 sensors for $59\%$, 10 sensors for $89\%$, and 22 sensors for $92\%$—higher density due to limited rooftop sites and deeper urban canyons.
\item \textbf{Height Placement}: Elevating sensors above average building heights (20-30\,m) remains the most effective lever for increasing $M_c^{20}$ without additional hardware. The altitude gap amplifies in real-world scenarios: Paulista's towers above 120\,m create 23.8 percentage-point differentials between 20\,m and 65\,m airways.
\item \textbf{Heterogeneous Mix}: Combining RF reach with Acoustic density and selective EO achieves higher redundancy than homogeneous deployments; Radar can remain optional for the scenarios analyzed. Real-world high-coverage solutions favor RF dominance (17 of 22 sensors) with supporting Radar, Acoustic, and EO nodes.
\item \textbf{Cost Target}: Operators should budget roughly \$4.4\,K per added coverage point up to $89\%$ (synthetic), and \$24\,K per point thereafter. Real-world dense corridors exhibit \$102\,K/pt premium for marginal gains above $89\%$, a 23$\times$ escalation that underscores the practical limits of sensor saturation in constrained urban environments.
\end{enumerate}

